%%%%%%%%%%%%%%%%%%%%%%%%%%%%%%%%%%%%%%%%%%%%%%%%%%%%%%%%%%%%%%%%%%%%%%
%                                                                    %
% Rational values of the arccosine function,                         %
% by                                                                 %
% J. L. Varona                                                       %
%                                                                    %
% Central European Journal of Mathematics,                           %
% Vol. 4 (2006), no. 2, 319-322.                                     %
%                                                                    %
% DATE OF LAST MODIFICATION: 13-January-2006                         %
% Written with LaTeX2e                                               %
%                                                                    %
%%%%%%%%%%%%%%%%%%%%%%%%%%%%%%%%%%%%%%%%%%%%%%%%%%%%%%%%%%%%%%%%%%%%%%

\documentclass[10pt,a4paper]{article}

\usepackage{amsmath,amsthm,amsfonts}

\newcommand{\R}{\mathbb{R}}
\newcommand{\Q}{\mathbb{Q}}
\newcommand{\Z}{\mathbb{Z}}
\newcommand{\N}{\mathbb{N}}

\newtheorem*{theorem}{Theorem}
\theoremstyle{remark}
\newtheorem*{rem}{Remark}

%	%%%%%%%%%%%%%%%%%%%%%%%%%%%%%%%%%%%%%%%%%%%%%%%%%%%%%%
\title{\bfseries Rational values of the arccosine function\thanks{Research partially 
supported by grant BFM2003-06335-C03-03 of the DGI (Spain).}}
%	% \title{Rational values of the arccosine function}
%	%%%%%%%%%%%%%%%%%%%%%%%%%%%%%%%%%%%%%%%%%%%%%%%%%%%%%%

%% WITH article:
\author{{\Large Juan L. Varona\thanks{Email: \texttt{jvarona@dmc.unirioja.es}}}\\[5pt]
\textit{\small Departamento de Matem\'aticas y Computaci\'on,}\\[-2pt]
\textit{\small Universidad de La Rioja,}\\[-2pt]
\textit{\small Edificio Vives, Calle Luis de Ulloa s/n,
26004 Logro\~no, Spain}}

\date{\small \leavevmode\\[5pt] \hfill Received 7 November 2005; 
accepted 8 February 2006}

%%%%%%%%%%%%%%%%%%%%%%%%%%%%%%%%%%%%%%%%%%%%%%%%%%%%%%
\begin{document}
%%%%%%%%%%%%%%%%%%%%%%%%%%%%%%%%%%%%%%%%%%%%%%%%%%%%%%

%%%%%%%%%%%%%%%%%%%%%%%%%%%%%%%%%%%%%%%%%%%%%%%%%%%%%%
\maketitle
%%%%%%%%%%%%%%%%%%%%%%%%%%%%%%%%%%%%%%%%%%%%%%%%%%%%%%

%%%%%%%%%%%%%%%%%%%%%%%%%%%%%%%%%%%%%%%%%%%%%%%%%%%%%%
\renewcommand\thefootnote{}
\renewcommand\footnotemark{}
\footnote{This paper has been published in: \textit{Cent.\@ Eur.\@ J.\@ Math.\@}
\textbf{4} (2006), no.~2, 319--322.}
%%%%%%%%%%%%%%%%%%%%%%%%%%%%%%%%%%%%%%%%%%%%%%%%%%%%%%

%%%%%%%%%%%%%%%%%%%%%%%%%%%%%%%%%%%%%%%%%%%%%%%%%%%%%%
\vspace*{-1cm}\hrule
\begin{abstract}\parindent=0pt\noindent
We give a short proof to characterize the cases when
$\arccos(\sqrt{r})$, the arccosine
of the square-root of a rational number~$r \in [0,1]$,
is a rational multiple of~$\pi$:
This happens exactly if $r$ is an integer multiple of~$1/4$.
The proof relies on the well-known recurrence relation for the
Chebyshev polynomials of the first kind.

\copyright\ Versita Warsaw and Springer-Verlag Berlin Heidelberg.
All rights reserved.
\bigskip

\textit{Keywords: Arccosine, cosine, rational, irrational.}

\textit{MSC (2000): 11J72, 33B10.}
\end{abstract}
\hrule\bigskip
\bigskip
%%%%%%%%%%%%%%%%%%%%%%%%%%%%%%%%%%%%%%%%%%%%%%%%%%%%%%

The arithmetic properties of trigonometric functions has been
a recurring topic in the mathematical literature. In 1933,
D.~H.~Lehmer \cite{Leh} proved that if $d>2$ and $k/d$ is an
irreducible fraction, then $2\cos(2\pi k/d)$ is an algebraic
integer of degree $\varphi(d)/2$ (with $\varphi(d)$ being
Euler's totient function); the proof can be also found
in~\cite[Theorem~3.9]{Ni}. As a consequence, it can be shown
that, for $t \in \Q$, the only rational values of
$\cos(\pi t)$ are $\cos(\pi t) = 0, \pm1, \pm1/2$. But this can
be proved independently of Lehmer's result; see also~\cite[Chapter~3]{Ni}
for historical references.
Another nice and self-contained proof appears
in~\cite[\S\,6.3, Theorem~6.16]{NiZuMo}.
% From this result and the relation $\cos^2(\theta) = (1+\cos(2\theta))/2$,
% we can describe the possible rational values of $\cos^2(\pi r)$
% when $r \in \Q$.

In~\cite[Chapter~6]{AiZi}, as a key step for the construction of Dehn's
counterexamples to Hilbert's third problem about decomposing 
polyhedra, it is established that
\begin{equation}
\label{eq:AiZi}
  \frac{1}{\pi} \arccos\left(\frac{1}{\sqrt{n}}\right) \notin \Q
  \quad\mbox{when }n \in \N,\quad n\mbox{ odd},\quad n \ge 3.
\end{equation}
%% NO ORIGINAL REFERENCE IS GIVEN IN~\cite[Chapter~6]{AiZi}
The aim of this paper is to give a direct and simple proof of
a much more general result: the complete characterization of the
$r \in \Q$ such that
$$
  \frac{1}{\pi} \arccos\left(\sqrt{r}\right) \in \Q.
$$

Let us explain the idea of our proof.
The elegant proof of~(\ref{eq:AiZi}) given in~\cite[Chapter~6]{AiZi}
is based in the trigonometric identity
\begin{equation}
\label{eq:coscos}
  \cos((k+1)\theta) = 2\cos(\theta)\cos(k\theta) - \cos((k-1)\theta),
\end{equation}
which is an immediate consequence of
$$
  \cos(\alpha) + \cos(\beta)
  = 2 \cos(\tfrac{\alpha+\beta}{2}) \cos(\tfrac{\alpha-\beta}{2}).
$$
For even~$n$, a different method is suggested in that book, distinguishing
between the cases $n=2^j$ and $n$ not a power of~$2$.
Thus, it is obtained that $(1/\pi)\arccos(1/\sqrt{n})$ is rational
if and only if $n \in \{1,2,4\}$.

The relation~(\ref{eq:coscos}) can be read in term of
Chebyshev polynomials of the first kind.
These polynomials are well known in the mathematical literature
(see, for instance, \cite{Wei} or~\cite{Te}),
mainly for their importance in approximation theory (they are orthogonal polynomials,
and are also used in least squares fit and in quatrature formulas for
numerical integration).
At least for this author, polynomial relations
are easier to handle than trigonometric relations and, thus, the proof
given in~\cite[Chapter~6]{AiZi} seems to be clearer when written in terms
of polynomials. Moreover, this allows a useful generalization
of~(\ref{eq:AiZi}) whose proof does not lose the simplicity of~\cite{AiZi};
concretely, in this way we show that
$$
  \frac{1}{\pi} \arccos\left(\frac{m}{2\sqrt{nM}}\right) \notin \Q
$$
when $n,m,M \in \N$, $n \ge 2$, $\gcd(n,m) = 1$ and $m/(2\sqrt{nM}) < 1$.

Finally, the square root of every positive rational number $r<1$
can be written as $\sqrt{r} = m/(2\sqrt{nM})$ with $\gcd(n,m) = 1$, with the exceptions
of $\sqrt{r} = 1$, $1/2$, $1/\sqrt{2}$ and $\sqrt{3}/2$. Thus, we conclude
that, for $r \in \Q$ with $0 \le r \le 1$, the number $(1/\pi)\arccos(\sqrt{r})$
is irrational except in the cases arising from these values of the cosine function:
% the following cases:
% the cases provided for these values of the cosine function:
\begin{equation}
\label{eq:cos}
\begin{aligned}
  &\cos(0) = 1,\quad \cos(\pi/6) = \sqrt{3}/2,\quad
  \cos(\pi/4) = 1/\sqrt{2},\\
  &\cos(\pi/3) = 1/2, \quad\mbox{and}\quad \cos(\pi/2) = 0.
\end{aligned}
\end{equation}

\begin{rem}
As an easy consequence, for $t \in \Q$, the only possible rational values of
$\cos^2(\pi t)$ are given by $\cos(\pi t) = \pm1$, $\pm\sqrt{3}/2$,
$\pm1/\sqrt{2}$, $\pm1/2$ and~$0$.
Actually, this can also be proved by using that the only rational values of
$\cos(\pi t)$ are $0$, $\pm1$, and $\pm1/2$, and the relation
$\cos^2(\theta) = (1+\cos(2\theta))/2$;
or can be derived from Lehmer's result by searching 
algebraic integers of degree at most~$2$ of $\cos(\pi t)$.
But it seems that none of these facts about
$\cos^2$ have been noticed in the literature. In any case,
we are giving a direct and new proof of this result.
Of course, using the elementary trigonometric relations
$\cos^2(\theta) = 1-\sin^2(\theta)$
and $\cos^2(\theta) = 1/(1 +\tan^2(\theta))$,
similar results for the rational values of $\sin^2(\pi t)$ and $\tan^2(\pi t)$
can be obtained.
\end{rem}

Thus, let us state

\begin{theorem}
Let $r \in \Q$ with $0 \le r \le 1$. Then, the number
$$
  \frac{1}{\pi} \arccos\left(\sqrt{r}\right)
$$
is rational if and only if $r$ is $0$, $1/4$, $1/2$, $3/4$, or~$1$;
%$$
%  \sqrt{r} \in \left\{ 0,\; 1,\; \frac12,\; \frac1{\sqrt{2}},\;
%  \frac3{2\sqrt{3}} \right\};
%$$
and the same holds for $(1/\pi) \arcsin(\sqrt{r})$.
\end{theorem}

\begin{proof}
Noticing that $\arccos(x) + \arcsin(x) = \pi/2$, it is enough to analyze
the case $\arccos$. Even more, the~``if'' part is clear from the trigonometric
values shown in~(\ref{eq:cos}), so let us study the ``only~if'' part.

We claim that every $r \in \Q \setminus \{0, 1/4, 2/4, 3/4, 1 \}$, $r \ge 0$,
can be written as
\begin{equation}
\label{eq:4nM}
  r = \frac{m^2}{4nM}
\end{equation}
with conditions $n,m,M \in \N$, $n \ge 2$ and $\gcd(n,m) = 1$.
This is true because given $r = p/q$, with $p$ and $q$ co-prime
and $q$ not a divisor of~$4$, there are only two possibilities:
\begin{itemize}
\item if $q$ has an odd divisor $n$, say $q = nM'$, we can write
$r = (2p)^2/(4nM)$ (with $M=pM'$);
\item if $q$ is a power of $2$, then $q = 2^j$ with $j \ge 3$,
$p$ is odd and we have $r = p^2/(4\cdot2^{j-2}p)$ (with $n = 2^{j-2} \ge 2$).
\end{itemize}
In both cases we have found the decomposition~(\ref{eq:4nM}).
Then, we only need to prove that
\begin{equation*}
\label{eq:nmM}
  A(n,m,M) = \frac{1}{\pi} \arccos\left(\frac{m}{2\sqrt{nM}}\right) \notin \Q
\end{equation*}
when $n,m,M \in \N$, $n \ge 2$, $\gcd(n,m) = 1$ and $m/(2\sqrt{nM}) < 1$.

For $x \in [-1,1]$ and $k \in \N \cup \{0\}$, let $T_k(x)$ be
$$
  T_k(\cos(\theta)) = \cos(k \theta), \qquad x = \cos(\theta).
$$
It is immediate that $T_0(x) = 1$ and $T_1(x) = x$. Moreover, the trigonometric relation
\begin{equation*}
  \cos((k+1)\theta) = 2\cos(\theta)\cos(k\theta) - \cos((k-1)\theta)
\end{equation*}
proves the recurrence formula
\begin{equation*}
%\label{eq:rec}
  T_{k+1}(x) = 2x T_k(x) - T_{k-1}(x), \quad k \ge 1.
\end{equation*}
In particular, this implies that $T_k(x)$ is a polynomial of degree~$k$,
and so we get the so-called Chebyshev polynomials of the first kind.
Making the substitution $g_k(x) = 2T_k(mx/2)$, we get
\begin{equation}
\label{eq:g_k}
\left\{
\begin{aligned}
  &g_0(x) = 2,\quad g_1(x) = m x,\\
  &g_{k+1}(x) = m x g_k(x) - g_{k-1}(x), \quad k \ge 1.
\end{aligned}
\right.
\end{equation}
Then, $g_k(x)$ is a polynomial of degree~$k$ and coefficients in~$\Z$,
and it verifies
\begin{equation}
\label{eq:2cosg}
  2 \cos(k \theta) = 2 T_k(\cos(\theta))
  = g_k(2\cos(\theta)/m).
\end{equation}

Now, let us take
$$
  \theta = \arccos\left(\frac{m}{2\sqrt{nM}}\right),
  \quad x = \frac{1}{2\sqrt{nM}},
  \quad \cos(\theta) = mx.
$$
By~(\ref{eq:2cosg}), we have
%% (remember that $g_k(x)$ is a polynomial of degree~$k$ and integer coefficients)
\begin{equation}
\label{eq:Bk}
  2 \cos(k \theta) = g_k(2x) = g_k\left(\frac{1}{\sqrt{nM}}\right)
  = \frac{B_k}{(\sqrt{nM})^k}
\end{equation}
for some~$B_k \in \Z$. From~(\ref{eq:g_k}) it is easy to check that
$B_0 = 2$, $B_1 = m$, and $B_{k+1} = m B_k - nM B_{k-1}$.
Now, let us recall that $\gcd(n,m) = 1$. Then, by induction on $k$ (starting with $k=1$),
it follows that $n$ does not divide $B_k$ for any $k \ge 1$.

To conclude the proof, let us suppose that $A(n,m,M) = (1/\pi) \theta = h/k \in \Q$.
Then, $k\theta = h\pi$ and, by~(\ref{eq:Bk}),
$$
  \pm 2 = 2 \cos(h\pi) = 2 \cos(k\theta) = \frac{B_k}{(\sqrt{nM})^k},
  \quad B_k \in \Z.
$$
This implies that $n$ divides $B_k$, which is a contradiction.
\end{proof}

\subsection*{Acknowledgment} Thanks to Professors G\"unter M.~Ziegler,
Jai\-me Vi\-nue\-sa, and the referee, for their
interest and their valuable comments, that allowed to improve this paper.
%% 
%% In ``{J. Vinuesa}, {Problems and solutions: Solution to the problem 28}
%% (Spanish), \textit{Gac. R. Soc. Mat. Esp.} \textbf{9} (2006), no.~1''
%% there appears the proof of the claim~(\ref{eq:4nM}). It was the solution
%% to the cooresponding problem proposed in ``{J. L. Varona}, {Problems and
%% solutions: Problem 28} (Spanish), \textit{Gac. R. Soc. Mat. Esp.}
%% \textbf{8} (2007), no.~1, p.~182''.
%% I tried to include the reference to the solucion of J. Vinuesa in the
%% final version of this paper, but unfortunately it were not possible.

%%%%%%%%%%%%%%%%%%%%%%%%%%%%%%%%%%%%%%%%%%%%%%%%%%%%%%
\begin{thebibliography}{9}
%%%%%%%%%%%%%%%%%%%%%%%%%%%%%%%%%%%%%%%%%%%%%%%%%%%%%%

\bibitem{AiZi}
{M. Aigner and G. M. Ziegler},
\textit{Proofs from THE BOOK},
3\textsuperscript{rd} edition, Springer, 2004.

\bibitem{Leh}
{D. H. Lehmer},
{A note on trigonometric algebraic numbers},
\textit{Amer. Math. Monthly}
\textbf{40} (1933), 165--166.

\bibitem{Ni}
{I. Niven},
\textit{Irrational numbers},
Carus Monographs no.~11,
The Mathematical Association of America
(distributed by John Wiley and Sons), 1956.

\bibitem{NiZuMo}
{I. Niven, H. S. Zuckerman, and H. L. Montgomery},
\textit{An introduction to the theory of numbers},
5\textsuperscript{th}~edition, Wiley, 1991.

\bibitem{Te}
{N. M. Temme},
\textit{Special functions: an introduction to the classical functions of
mathematical physics},
John Wiley and Sons, 1996.

\bibitem{Wei}
{E. W. Weisstein},
{Chebyshev polynomial of the first kind}.
From \textit{MathWorld}--A Wolfram Web Resource,
\texttt{http://mathworld.wolfram.com/\allowbreak ChebyshevPolynomialoftheFirstKind.html}.

%%%%%%%%%%%%%%%%%%%%%%%%%%%%%%%%%%%%%%%%%%%%%%%%%%%%%%
\end{thebibliography}
%%%%%%%%%%%%%%%%%%%%%%%%%%%%%%%%%%%%%%%%%%%%%%%%%%%%%%

%%%%%%%%%%%%%%%%%%%%%%%%%%%%%%%%%%%%%%%%%%%%%%%%%%%%%%
\end{document}
%%%%%%%%%%%%%%%%%%%%%%%%%%%%%%%%%%%%%%%%%%%%%%%%%%%%%%


